\begin{textsample}{POS Dim 3 – rj – Score 18.00 – rj/rj\_inf\_camp\_exp\_8.txt}  \label{ex:f3_pos_008}
O Campo de Experiências Traços, \textbf{Sons}, Cores e Formas A BNCC, assim apresenta esse campo : > Conviver com diferentes manifestações artísticas, culturais e científicas, locais e universais, no cotidiano da instituição escolar, possibilita às \textbf{crianças}, por meio de experiências diversificadas, vivenciar diversas formas de expressão e linguagens, como as artes visuais ( \textbf{pintura}, modelagem, colagem, fotografia etc. ), a música, o teatro, a dança e o audiovisual, entre outras.
Com base nessas experiências, elas se expressam por várias linguagens, criando suas próprias produções artísticas ou culturais, exercitando a autoria ( coletiva e individual ) com \textbf{sons}, traços, \textbf{gestos}, danças, \textbf{mímicas}, encenações, canções, desenhos, modelagens, manipulação de diversos materiais e de recursos tecnológicos.
Essas experiências contribuem para que, desde muito \textbf{pequenas}, as \textbf{crianças} desenvolvem senso estético e crítico, o conhecimento de si mesmas, dos outros e da realidade que as cerca.
Portanto a Educação \textbf{Infantil} precisa promover a participação das \textbf{crianças} em tempos e espaços para a produção, manifestação e apreciação artística, de modo a favorecer o desenvolvimento da sensibilidade, da criatividade e da expressão pessoal das \textbf{crianças}, permitindo que elas se apropriem e reconfigurem, permanentemente, a cultura e potencializem suas singularidades, ao ampliar repertórios e interpretar suas experiências artísticas.
As \textbf{crianças} vivem em ambientes onde a cada momento ocorrem situações envolvendo pessoas, atividades, espaços, objetos e materiais que elas buscam perceber, reconhecer, significar e representar, e o fazem pela apropriação de diferentes linguagens e recursos, como suas sensações, \textbf{afetos} e \textbf{desejos}, sua corporeidade, sua linguagem verbal, sua percepção das ações de seus parceiros e sua atenção voltada para os aspectos materiais do ambiente.
Essas aquisições pela \textbf{criança} dependem em especial das interações que ela estabelece com seus parceiros, dado que as características das atividades, das ações dos que as executam e dos aspectos físicos e simbólicos dos espaços em que elas ocorrem, levam a \textbf{criança} a significar este ambiente como um local que amedronta, ou que é divertido, conforme sinta que foi afetada por ele.
Dentre outros elementos, o ambiente que a \textbf{criança} percebe é composto de traços, \textbf{sons}, cores e formas, que aparecem também nas plantas, solos e outros elementos do entorno, em objetos e \textbf{vestimentas}, em gravuras de livros, na caracterização própria de atores e cenários teatrais, na decoração e iluminação do espaço etc.
O ambiente envolve a \textbf{criança} em um mundo de estímulos visuais, sonoros, táteis, cinéticos em que o formato, a cor, a \textbf{textura}, a sonoridade e o uso habitual dos objetos – \textbf{brinquedos}, latas, blocos de madeira, telas, argila, livros, cartazes, \textbf{lápis}, canetas, pedaços de giz ou carvão –, a presença de equipamentos elétricos e eletrônicos, e de \textbf{mobiliário} e adornos presentes na ocorrência de certas atividades, provocam nela emoções e representações.
Perceber os objetos a partir de sua estrutura, de sua forma, de suas qualidades \textbf{sensoriais}, e não apenas segundo sua utilização imediata, possibilita às \textbf{crianças} ver o mundo como uma soma de estímulos e não como mera série de \textbf{utensílios}, dentro de um projeto pedagógico que não é voltado à produção de apresentações em datas comemorativas, ou de “ trabalhinhos ” a serem levados para casa.
Essa visão supera tendências de apresentar as \textbf{crianças} tão somente o mundo das utilidades, da produção industrial, e as \textbf{apoia} a pensar o mundo como um espaço para ver o ambiente com sensibilidade, como uma paisagem cativante, fugindo de padrões estereotipados e de elementos da cultura de massa.

% matched lemmas: afeto, apoiar, brinquedo, criança, desejo, gesto, infantil, lápis, mobiliário, mímico, pequeno, pintura, sensorial, som, textura, utensílio, vestimenta
\end{textsample}
